\documentclass[12pt]{article}
\usepackage[utf8]{inputenc}
\usepackage[margin=1in]{geometry}
\usepackage{amsmath, amssymb, amsthm}
\usepackage{xcolor}
\usepackage[most]{tcolorbox}
\usepackage{enumitem}
\usepackage{fancyhdr}
\usepackage{titlesec}

% Define colored boxes matching the framework
\newtcolorbox{problemconfigbox}[1]{
    colback=blue!5!white,
    colframe=blue!75!black,
    title=#1,
    fonttitle=\bfseries,
    breakable
}

\newtcolorbox{strategicbox}[1]{
    colback=pink!10!white,
    colframe=pink!75!black,
    title=#1,
    fonttitle=\bfseries,
    breakable
}

\newtcolorbox{tacticalbox}[1]{
    colback=green!5!white,
    colframe=green!75!black,
    title=#1,
    fonttitle=\bfseries,
    breakable
}

\newtcolorbox{conversionbox}[1]{
    colback=yellow!10!white,
    colframe=yellow!75!black,
    title=#1,
    fonttitle=\bfseries,
    breakable
}

\newtcolorbox{verificationbox}[1]{
    colback=red!5!white,
    colframe=red!75!black,
    title=#1,
    fonttitle=\bfseries,
    breakable
}

\newtcolorbox{roadmapbox}[1]{
    colback=cyan!5!white,
    colframe=cyan!75!black,
    title=#1,
    fonttitle=\bfseries,
    breakable
}

\newtcolorbox{competitionbox}[1]{
    colback=purple!5!white,
    colframe=purple!75!black,
    title=#1,
    fonttitle=\bfseries,
    breakable
}

\newtcolorbox{insightbox}{
    colback=orange!5!white,
    colframe=orange!75!black,
    title=\textbf{KEY INSIGHT},
    fonttitle=\bfseries,
    breakable
}

\newtcolorbox{warningbox}{
    colback=red!10!white,
    colframe=red!75!black,
    title=\textbf{WARNING},
    fonttitle=\bfseries,
    breakable
}

\newtcolorbox{protipbox}{
    colback=teal!5!white,
    colframe=teal!75!black,
    title=\textbf{PRO TIP},
    fonttitle=\bfseries,
    breakable
}

\title{\textbf{AMC12A 2016 Problem 16:\\Strategic Analysis Using Comprehensive Framework}}
\author{Problem-Solving Strategic Analysis}
\date{\today}

\begin{document}

\maketitle

\section*{Problem Statement}

The graphs of $y=\log_3 x$, $y=\log_x 3$, $y=\log_{\frac{1}{3}} x$, and $y=\log_x \frac{1}{3}$ are plotted on the same set of axes. How many points in the plane with positive $x$-coordinates lie on two or more of the graphs?

$\textbf{(A)}\ 2\qquad\textbf{(B)}\ 3\qquad\textbf{(C)}\ 4\qquad\textbf{(D)}\ 5\qquad\textbf{(E)}\ 6$

\newpage

\section{PROBLEM CONFIGURATION ANALYSIS}

\begin{problemconfigbox}{A. Problem Classification}
\begin{itemize}[leftmargin=*]
    \item \textbf{Problem Type}: To Find - counting the number of intersection points among four logarithmic graphs
    
    \item \textbf{Difficulty Band}: Late-band (Problem 16 out of 25), requiring sophisticated understanding of logarithmic properties and systematic case analysis
    
    \item \textbf{Competition Context}: AMC12 (75 minutes, 25 problems) - approximately 3 minutes available for this problem, but as a P16, likely requires 4-6 minutes for most students
    
    \item \textbf{Mathematical Domain}: Algebra (logarithms), with elements of function analysis and equation solving
    
    \item \textbf{Source Context}: 2016 AMC12A - historically, P16 has 15-25\% solve rate among all participants
\end{itemize}
\end{problemconfigbox}

\begin{problemconfigbox}{B. Problem Structure Identification}
\begin{itemize}[leftmargin=*]
    \item \textbf{Given Information}: 
    \begin{itemize}
        \item Four functions: $f_1(x) = \log_3 x$, $f_2(x) = \log_x 3$, $f_3(x) = \log_{\frac{1}{3}} x$, $f_4(x) = \log_x \frac{1}{3}$
        \item All graphs plotted on same axes
        \item Domain restriction: $x > 0$ (positive $x$-coordinates)
        \item Additional implicit restriction: $x \neq 1$ (logarithms with base 1 undefined)
    \end{itemize}
    
    \item \textbf{Constraints}: 
    \begin{itemize}
        \item $x > 0$ and $x \neq 1$ for all functions
        \item For $f_2$ and $f_4$: base is $x$, so $x > 0$ and $x \neq 1$
        \item Must find intersections, not just points on individual graphs
    \end{itemize}
    
    \item \textbf{Objective}: Count distinct points $(x, y)$ where at least two of the four graphs intersect, considering only positive $x$-coordinates
    
    \item \textbf{Hidden Assumptions}: 
    \begin{itemize}
        \item Each pair of functions can intersect at most at a finite number of points
        \item Point $(1, 0)$ might seem like a candidate but $x = 1$ is outside domain for $f_2$ and $f_4$
        \item Some intersections might coincide (same point for multiple pairs)
    \end{itemize}
    
    \item \textbf{Answer Choice Leverage}: 
    \begin{itemize}
        \item Choices range from 2 to 6, suggesting moderate number of intersections
        \item Can verify answer by systematic enumeration
        \item If computed answer doesn't match choices, signals calculation error
    \end{itemize}
\end{itemize}
\end{problemconfigbox}

\begin{problemconfigbox}{C. Strategic Orientation}
\begin{itemize}[leftmargin=*]
    \item \textbf{Hypothesis}: We need to check all $\binom{4}{2} = 6$ pairs of functions for intersections
    
    \item \textbf{Conclusion}: Find the total count of distinct intersection points (with positive $x$-coordinates)
    
    \item \textbf{Notation}: 
    \begin{itemize}
        \item Let $f_1(x) = \log_3 x$, $f_2(x) = \log_x 3$, $f_3(x) = \log_{\frac{1}{3}} x$, $f_4(x) = \log_x \frac{1}{3}$
        \item For any intersection: $f_i(x) = f_j(x)$ for some $i \neq j$
    \end{itemize}
    
    \item \textbf{Intuitive Approach}: 
    \begin{itemize}
        \item Use logarithm change-of-base formula and properties
        \item Look for relationships between the four functions
        \item Systematically check all pairs
        \item Watch for symmetries that might reduce work
    \end{itemize}
    
    \item \textbf{Cross-Domain Potential}: 
    \begin{itemize}
        \item Could transform to exponential equations
        \item Could use substitution $t = \log_3 x$ to simplify
        \item Graphical visualization might reveal patterns
    \end{itemize}
\end{itemize}
\end{problemconfigbox}

\section{STRATEGIC FRAMEWORK APPLICATION}

\begin{strategicbox}{A. Psychological Strategies}
\begin{itemize}[leftmargin=*]
    \item \textbf{Mental Toughness}: This problem requires checking 6 pairs of functions. Don't give up after checking 2-3 pairs. Systematic enumeration is key. If initial approach seems tedious, look for relationships between functions to reduce work.
    
    \item \textbf{Creativity Cultivation}: 
    \begin{itemize}
        \item Notice that $\log_{\frac{1}{3}} x = -\log_3 x$ (reciprocal base property)
        \item Notice that $\log_x 3 = \frac{1}{\log_3 x}$ (change of base)
        \item Look for function families or symmetries
        \item Consider making substitution $t = \log_3 x$ to simplify all four functions
    \end{itemize}
    
    \item \textbf{Iterative Refinement}: Start with easier pairs (functions with same base structure), gain insight, then tackle harder pairs. If algebra becomes messy, verify numerically with test values.
\end{itemize}
\end{strategicbox}

\begin{strategicbox}{B. Investigation Strategies}
\begin{itemize}[leftmargin=*]
    \item \textbf{Startup Orientation}: Read problem carefully. Identify that we need intersection points, not just points on graphs. Note the restriction $x > 0$.
    
    \item \textbf{Get Your Hands Dirty}: 
    \begin{itemize}
        \item Test $x = 3$: $\log_3 3 = 1$, $\log_3 3 = 1$, $\log_{\frac{1}{3}} 3 = -1$, $\log_3 \frac{1}{3} = -1$
        \item Test $x = \frac{1}{3}$: $\log_3 \frac{1}{3} = -1$, $\log_{\frac{1}{3}} 3 = 1$, $\log_{\frac{1}{3}} \frac{1}{3} = 1$, $\log_{\frac{1}{3}} \frac{1}{3} = -1$
        \item Already finding patterns!
    \end{itemize}
    
    \item \textbf{Penultimate Step}: Before counting total intersections, we need to find all intersection points. The penultimate step is solving each of the 6 pairwise equations.
    
    \item \textbf{Wishful Thinking}: Wish that logarithmic properties simplify the four functions into a smaller set of related expressions (they do!)
    
    \item \textbf{Make It Easier}: 
    \begin{itemize}
        \item Use substitution $t = \log_3 x$ to convert all functions
        \item This transforms: $f_1 \to y = t$, $f_2 \to y = \frac{1}{t}$, $f_3 \to y = -t$, $f_4 \to y = -\frac{1}{t}$
        \item Now problem becomes: find intersections of $y = t$, $y = \frac{1}{t}$, $y = -t$, $y = -\frac{1}{t}$
    \end{itemize}
    
    \item \textbf{Conjecture via Small Cases}: From test values, conjecture that $x = 3$ and $x = \frac{1}{3}$ are intersection points. Verify systematically.
    
    \item \textbf{Visualization}: Sketch the four curves. Notice that $f_1$ and $f_3$ are reflections across $x$-axis, while $f_2$ and $f_4$ are reflections across $x$-axis of each other. Also $f_1$ and $f_2$ are reflections across $y = x$.
\end{itemize}
\end{strategicbox}

\begin{strategicbox}{C. Late-Band Strategic Playbook}
\begin{itemize}[leftmargin=*]
    \item \textbf{Answer-Choice Leverage}: 
    \begin{itemize}
        \item If we find 4 intersections quickly but answer choices include 5, we know to look harder
        \item Extreme answers (2 or 6) would suggest either many cancellations or all pairs intersect
        \item Middle answer (4) suggests moderate complexity
    \end{itemize}
    
    \item \textbf{Make It Easier}: Transform to $t$-coordinates as described above. This dramatically simplifies the algebra.
    
    \item \textbf{Penultimate Step}: The step before counting is ensuring we haven't double-counted any points and that we've checked all 6 pairs.
    
    \item \textbf{Wishful Thinking}: Assume the functions have nice relationships (which they do via reciprocal and negative relationships).
    
    \item \textbf{Conjecture via Small Cases}: Testing $x = 3, \frac{1}{3}, 1$ reveals most intersections immediately.
\end{itemize}
\end{strategicbox}

\begin{strategicbox}{D. Argument Construction Strategy}
\begin{itemize}[leftmargin=*]
    \item \textbf{Direct Approach}: 
    \begin{enumerate}
        \item Use logarithm properties to relate all four functions
        \item Apply substitution $t = \log_3 x$ to simplify
        \item Systematically solve all 6 pairwise equations in $t$-coordinates
        \item Convert back to $x$-coordinates for each solution
        \item Verify each solution satisfies domain constraints
        \item Count distinct solutions
    \end{enumerate}
    
    \item \textbf{Systematic Enumeration}: Check all $\binom{4}{2} = 6$ pairs, being careful not to miss any or double-count.
\end{itemize}
\end{strategicbox}

\section{TACTICAL METHODOLOGY SELECTION}

\begin{tacticalbox}{A. Core Mathematical Tactics}
\begin{itemize}[leftmargin=*]
    \item \textbf{Symmetry Exploitation}:
    \begin{itemize}
        \item $\log_{\frac{1}{a}} b = -\log_a b$ (reciprocal base property)
        \item $\log_a b \cdot \log_b a = 1$ (change of base reciprocal property)
        \item These symmetries reduce the complexity of the problem
    \end{itemize}
    
    \item \textbf{Change of Variables}:
    \begin{itemize}
        \item Substitution $t = \log_3 x$ linearizes some functions
        \item Transforms domain: $x \in (0, 1) \cup (1, \infty) \to t \in \mathbb{R} \setminus \{0\}$
        \item Key transformations: $\log_x 3 = \frac{1}{t}$, $\log_{\frac{1}{3}} x = -t$, $\log_x \frac{1}{3} = -\frac{1}{t}$
    \end{itemize}
    
    \item \textbf{Systematic Casework}:
    \begin{itemize}
        \item Organize pairs: $(f_1, f_2)$, $(f_1, f_3)$, $(f_1, f_4)$, $(f_2, f_3)$, $(f_2, f_4)$, $(f_3, f_4)$
        \item Solve each equation separately
        \item Check that solutions don't coincide
    \end{itemize}
    
    \item \textbf{Algebraic Manipulation}:
    \begin{itemize}
        \item Use $\log_a b = \frac{\log b}{\log a}$ (change of base formula)
        \item Recognize when equations simplify to polynomial equations in $t$
    \end{itemize}
\end{itemize}
\end{tacticalbox}

\begin{tacticalbox}{B. Domain-Specific Techniques}
\begin{itemize}[leftmargin=*]
    \item \textbf{Logarithmic Properties}:
    \begin{itemize}
        \item $\log_a b = \frac{1}{\log_b a}$
        \item $\log_{\frac{1}{a}} b = -\log_a b$
        \item $\log_a a = 1$ and $\log_a 1 = 0$
    \end{itemize}
    
    \item \textbf{Function Analysis}:
    \begin{itemize}
        \item Identify increasing vs decreasing functions
        \item $f_1$ is increasing, $f_3$ is decreasing
        \item $f_2$ is decreasing on $(0, 1)$ and $(1, \infty)$, undefined at $x = 1$
        \item $f_4$ is increasing on $(0, 1)$ and $(1, \infty)$, undefined at $x = 1$
    \end{itemize}
\end{itemize}
\end{tacticalbox}

\section{CONVERSION TECHNIQUES APPLICATION}

\begin{conversionbox}{A. Recognition Index and Technique Selection}
\begin{itemize}[leftmargin=*]
    \item \textbf{High-Probability Techniques Applied}:
    \begin{enumerate}
        \item \textbf{Change of Base Formula} (95\% confidence): Essential for relating $\log_3 x$ and $\log_x 3$
        \item \textbf{Reciprocal Base Property} (95\% confidence): $\log_{\frac{1}{3}} x = -\log_3 x$ simplifies half the functions
        \item \textbf{Substitution Method} (90\% confidence): $t = \log_3 x$ reduces to simpler algebraic equations
        \item \textbf{Systematic Enumeration} (85\% confidence): Checking all 6 pairs ensures completeness
    \end{enumerate}
    
    \item \textbf{Medium-Probability Techniques}:
    \begin{enumerate}
        \item \textbf{Symmetry Analysis} (70\% confidence): Recognizing function relationships reduces computational burden
        \item \textbf{Graphical Visualization} (60\% confidence): Mental picture helps avoid missing intersections
    \end{enumerate}
    
    \item \textbf{Why These Techniques Apply}:
    \begin{itemize}
        \item Problem explicitly involves logarithms with different bases → change of base essential
        \item Reciprocal bases $3$ and $\frac{1}{3}$ → reciprocal base property naturally applies
        \item Need to find all intersections → systematic enumeration prevents missing cases
        \item Multiple functions with common structure → substitution simplifies
    \end{itemize}
\end{itemize}
\end{conversionbox}

\begin{conversionbox}{B. Technique Integration Strategy}
\begin{itemize}[leftmargin=*]
    \item \textbf{Primary Technique}: Substitution $t = \log_3 x$ combined with logarithmic properties
    
    \item \textbf{Integration Sequence}:
    \begin{enumerate}
        \item Apply logarithmic properties to express all functions in terms of $\log_3 x$
        \item Introduce substitution $t = \log_3 x$ to get: $y = t$, $y = \frac{1}{t}$, $y = -t$, $y = -\frac{1}{t}$
        \item Solve pairwise equations in $t$-$y$ space
        \item Convert solutions back to $x$-$y$ space using $x = 3^t$
        \item Verify all solutions have $x > 0$ and $x \neq 1$
    \end{enumerate}
    
    \item \textbf{Rationale for Integration}:
    \begin{itemize}
        \item Substitution converts transcendental equations to algebraic equations
        \item Working in $t$-space makes symmetries more apparent
        \item Systematic enumeration in simplified space is more reliable
    \end{itemize}
\end{itemize}
\end{conversionbox}

\begin{conversionbox}{C. Conflict Resolution}
\begin{itemize}[leftmargin=*]
    \item \textbf{Potential Conflict}: Whether to work in $x$-space or $t$-space
    \begin{itemize}
        \item \textbf{Resolution}: Work primarily in $t$-space for simplicity, convert final answers to $x$-space
        \item $t$-space equations are polynomial/rational, easier to solve
    \end{itemize}
    
    \item \textbf{Potential Conflict}: Domain restrictions when $x = 1$ (equivalently $t = 0$)
    \begin{itemize}
        \item \textbf{Resolution}: Explicitly exclude $t = 0$ from solutions since $x = 1$ makes $f_2$ and $f_4$ undefined
        \item Check this when solving $\frac{1}{t} = -t$ (gives $t^2 = -1$, no real solutions anyway)
    \end{itemize}
    
    \item \textbf{Potential Conflict}: Double-counting intersections
    \begin{itemize}
        \item \textbf{Resolution}: After finding all pairwise intersections, check if any points appear in multiple pairs
        \item Keep organized list of distinct $(x, y)$ coordinates
    \end{itemize}
\end{itemize}
\end{conversionbox}

\section{VERIFICATION STRATEGY DESIGN}

\begin{verificationbox}{A. Verification Level Selection}
\begin{itemize}[leftmargin=*]
    \item \textbf{Decision Tree Application}:
    \begin{itemize}
        \item Problem difficulty: Late-band (P16) → suggests Level 2 or Level 3
        \item Confidence after solution: If systematic and organized → Level 2
        \item Time remaining: If $<$ 20 minutes remain → Level 1
        \item Stakes: AMC12 P16 has moderate stakes → Level 2 appropriate
    \end{itemize}
    
    \item \textbf{Selected Level}: Level 2 - Structured Verification
    \begin{itemize}
        \item Check each pairwise solution by substitution
        \item Verify no solutions were double-counted
        \item Confirm all solutions satisfy $x > 0$ and $x \neq 1$
        \item Quick sanity check: does answer match an option?
    \end{itemize}
\end{itemize}
\end{verificationbox}

\begin{verificationbox}{B. Specialized Verification Methods}
\begin{itemize}[leftmargin=*]
    \item \textbf{Algebraic Verification}:
    \begin{itemize}
        \item For each intersection point $(x_0, y_0)$, verify that $y_0 = f_i(x_0) = f_j(x_0)$
        \item Example: For $(3, 1)$, check $\log_3 3 = 1$ and $\log_3 3 = 1$ $\checkmark$
    \end{itemize}
    
    \item \textbf{Completeness Check}:
    \begin{itemize}
        \item List all 6 pairs: $(f_1, f_2)$, $(f_1, f_3)$, $(f_1, f_4)$, $(f_2, f_3)$, $(f_2, f_4)$, $(f_3, f_4)$
        \item Confirm each pair was solved
        \item Note which pairs yield intersections and which don't
    \end{itemize}
    
    \item \textbf{Boundary Case Verification}:
    \begin{itemize}
        \item Check behavior as $x \to 0^+$: different functions diverge differently
        \item Check behavior as $x \to \infty$: again, different asymptotic behavior
        \item Confirm no intersections were missed at extremes
    \end{itemize}
    
    \item \textbf{Graphical Sanity Check}:
    \begin{itemize}
        \item Mental sketch: $f_1$ increasing through origin, $f_2$ hyperbola-like, $f_3$ decreasing, $f_4$ inverted hyperbola
        \item Do claimed intersections make geometric sense?
        \item Count should align with visual intuition
    \end{itemize}
\end{itemize}
\end{verificationbox}

\begin{verificationbox}{C. Confidence Calibration}
\begin{itemize}[leftmargin=*]
    \item \textbf{Initial Confidence}: 60\% (before solving)
    \begin{itemize}
        \item Logarithm problems can be tricky with domain restrictions
        \item Systematic enumeration should work but easy to make errors
    \end{itemize}
    
    \item \textbf{Post-Solution Confidence}: 85\% (after solving with substitution method)
    \begin{itemize}
        \item Substitution method greatly simplified the problem
        \item Found 5 distinct intersections systematically
        \item Answer matches option (D)
    \end{itemize}
    
    \item \textbf{Post-Verification Confidence}: 95\%
    \begin{itemize}
        \item All 5 points verified by direct substitution
        \item Confirmed all 6 pairs were checked
        \item No double-counting detected
        \item Domain restrictions satisfied for all solutions
    \end{itemize}
    
    \item \textbf{Decision}: Select answer (D) 5 with high confidence
\end{itemize}
\end{verificationbox}

\section{SOLUTION ROADMAP}

\begin{roadmapbox}{A. Detailed Step-by-Step Plan}
\begin{enumerate}[leftmargin=*]
    \item \textbf{Step 1: Apply Logarithmic Properties} (Time: 1 minute)
    \begin{itemize}
        \item Express all four functions in terms of $\log_3 x$:
        \item $f_1(x) = \log_3 x$
        \item $f_2(x) = \log_x 3 = \frac{1}{\log_3 x}$ (change of base)
        \item $f_3(x) = \log_{\frac{1}{3}} x = -\log_3 x$ (reciprocal base)
        \item $f_4(x) = \log_x \frac{1}{3} = -\frac{1}{\log_3 x}$ (combining both properties)
    \end{itemize}
    
    \item \textbf{Step 2: Introduce Substitution} (Time: 30 seconds)
    \begin{itemize}
        \item Let $t = \log_3 x$
        \item Then $x = 3^t$, and domain $x \in (0, 1) \cup (1, \infty)$ becomes $t \in \mathbb{R} \setminus \{0\}$
        \item Functions become:
        \begin{align*}
        y &= t \\
        y &= \frac{1}{t} \\
        y &= -t \\
        y &= -\frac{1}{t}
        \end{align*}
    \end{itemize}
    
    \item \textbf{Step 3: Solve Pairwise Intersections} (Time: 2-3 minutes)
    \begin{itemize}
        \item \textbf{Pair $(f_1, f_2)$}: $t = \frac{1}{t} \Rightarrow t^2 = 1 \Rightarrow t = \pm 1$
        \begin{itemize}
            \item $t = 1$: $x = 3^1 = 3$, $y = 1$ → point $(3, 1)$ $\checkmark$
            \item $t = -1$: $x = 3^{-1} = \frac{1}{3}$, $y = -1$ → point $(\frac{1}{3}, -1)$ $\checkmark$
        \end{itemize}
        
        \item \textbf{Pair $(f_1, f_3)$}: $t = -t \Rightarrow 2t = 0 \Rightarrow t = 0$
        \begin{itemize}
            \item But $t = 0$ means $x = 1$, which is outside domain (makes $f_2, f_4$ undefined)
            \item No valid intersection $\times$
        \end{itemize}
        
        \item \textbf{Pair $(f_1, f_4)$}: $t = -\frac{1}{t} \Rightarrow t^2 = -1$
        \begin{itemize}
            \item No real solutions $\times$
        \end{itemize}
        
        \item \textbf{Pair $(f_2, f_3)$}: $\frac{1}{t} = -t \Rightarrow 1 = -t^2 \Rightarrow t^2 = -1$
        \begin{itemize}
            \item No real solutions $\times$
        \end{itemize}
        
        \item \textbf{Pair $(f_2, f_4)$}: $\frac{1}{t} = -\frac{1}{t} \Rightarrow \frac{2}{t} = 0$
        \begin{itemize}
            \item No solutions $\times$
        \end{itemize}
        
        \item \textbf{Pair $(f_3, f_4)$}: $-t = -\frac{1}{t} \Rightarrow t = \frac{1}{t} \Rightarrow t^2 = 1 \Rightarrow t = \pm 1$
        \begin{itemize}
            \item $t = 1$: $x = 3$, $y = -1$ → point $(3, -1)$ $\checkmark$
            \item $t = -1$: $x = \frac{1}{3}$, $y = 1$ → point $(\frac{1}{3}, 1)$ $\checkmark$
        \end{itemize}
    \end{itemize}
    
    \item \textbf{Step 4: Check for Point at $x = 1$} (Time: 30 seconds)
    \begin{itemize}
        \item Could $(1, 0)$ be an intersection point?
        \item $f_1(1) = \log_3 1 = 0$ $\checkmark$
        \item $f_2(1) = \log_1 3$ is undefined (base cannot be 1) $\times$
        \item $f_3(1) = \log_{\frac{1}{3}} 1 = 0$ $\checkmark$
        \item $f_4(1) = \log_1 \frac{1}{3}$ is undefined $\times$
        \item So $f_1$ and $f_3$ intersect at $(1, 0)$, giving us our 5th point: $(1, 0)$ $\checkmark$
        \item Note: This intersection occurs at the boundary of the domain for $f_2$ and $f_4$
    \end{itemize}
    
    \item \textbf{Step 5: Compile Distinct Points} (Time: 30 seconds)
    \begin{itemize}
        \item From $(f_1, f_2)$: $(3, 1)$ and $(\frac{1}{3}, -1)$
        \item From $(f_1, f_3)$: $(1, 0)$
        \item From $(f_3, f_4)$: $(3, -1)$ and $(\frac{1}{3}, 1)$
        \item Total: 5 distinct points
        \item Verify no duplicates: all 5 points have different $(x, y)$ coordinates $\checkmark$
    \end{itemize}
    
    \item \textbf{Step 6: Verify Each Point} (Time: 1 minute)
    \begin{itemize}
        \item $(3, 1)$: $\log_3 3 = 1$ $\checkmark$ and $\log_3 3 = 1$ $\checkmark$ → $f_1$ and $f_2$ intersect
        \item $(\frac{1}{3}, -1)$: $\log_3 \frac{1}{3} = -1$ $\checkmark$ and $\log_{\frac{1}{3}} 3 = -1$ $\checkmark$ → $f_1$ and $f_2$ intersect
        \item $(1, 0)$: $\log_3 1 = 0$ $\checkmark$ and $\log_{\frac{1}{3}} 1 = 0$ $\checkmark$ → $f_1$ and $f_3$ intersect
        \item $(3, -1)$: $\log_{\frac{1}{3}} 3 = -1$ $\checkmark$ and $\log_3 \frac{1}{3} = -1$ $\checkmark$ → $f_3$ and $f_4$ intersect
        \item $(\frac{1}{3}, 1)$: $\log_{\frac{1}{3}} \frac{1}{3} = 1$ $\checkmark$ and $\log_{\frac{1}{3}} \frac{1}{3} = 1$ $\checkmark$ → $f_3$ and $f_4$ intersect
    \end{itemize}
\end{enumerate}

\textbf{Checkpoint}: After Step 3, should have 4 points. After Step 4, should have 5 points total.

\textbf{Time Management}: Total estimated time 5-6 minutes for complete solution.
\end{roadmapbox}

\begin{roadmapbox}{B. Alternative Approaches}
\begin{itemize}[leftmargin=*]
    \item \textbf{Alternative 1: Direct Algebraic Approach (without substitution)}
    \begin{itemize}
        \item Solve each pair directly: e.g., $\log_3 x = \log_x 3$
        \item Use change of base: $\log_3 x = \frac{1}{\log_3 x} \Rightarrow (\log_3 x)^2 = 1$
        \item More tedious but arrives at same answer
        \item Use if substitution method feels unfamiliar
    \end{itemize}
    
    \item \textbf{Alternative 2: Graphical Approach}
    \begin{itemize}
        \item Sketch all four functions on coordinate plane
        \item Visually identify intersection points
        \item Verify algebraically
        \item Good for visual learners but time-consuming
        \item Risk of missing intersections or miscounting
    \end{itemize}
    
    \item \textbf{Alternative 3: Answer Choice Elimination}
    \begin{itemize}
        \item Quickly find obvious intersections (e.g., at $x = 3$)
        \item Estimate if more exist
        \item Use symmetry arguments to predict total
        \item Only viable under extreme time pressure
    \end{itemize}
\end{itemize}
\end{roadmapbox}

\begin{roadmapbox}{C. Pivot Indicators and Backup Plans}
\begin{itemize}[leftmargin=*]
    \item \textbf{Pivot Indicator 1}: If after 2 minutes, only found 1-2 intersections
    \begin{itemize}
        \item \textbf{Action}: Switch to substitution method if using direct approach
        \item Or sketch graphs to gain visual insight
    \end{itemize}
    
    \item \textbf{Pivot Indicator 2}: If computed answer doesn't match any option
    \begin{itemize}
        \item \textbf{Action}: Check for double-counting or missed cases
        \item Verify domain restrictions were applied correctly
        \item Recount distinct points
    \end{itemize}
    
    \item \textbf{Pivot Indicator 3}: If algebra becomes too complex
    \begin{itemize}
        \item \textbf{Action}: Test numerical values (e.g., $x = 2, 3, 4, \frac{1}{2}, \frac{1}{3}$)
        \item Use numerical evidence to guide algebraic solving
    \end{itemize}
    
    \item \textbf{Abandonment Criterion}: If after 7 minutes no clear path to solution
    \begin{itemize}
        \item Use answer choice elimination with partial work
        \item Mark answer and move on
        \item Return if time permits at end
    \end{itemize}
\end{itemize}
\end{roadmapbox}

\begin{roadmapbox}{D. Comprehensive Pitfall Guide}
\begin{enumerate}[leftmargin=*]
    \item \textbf{Pitfall}: Forgetting domain restriction $x \neq 1$
    \begin{itemize}
        \item \textbf{Impact}: HIGH - Could include $(1, 0)$ incorrectly without checking $f_2$ and $f_4$
        \item \textbf{Prevention}: Always check that all functions are defined at claimed intersection
        \item \textbf{Detection}: Verify each point by substitution into all relevant functions
    \end{itemize}
    
    \item \textbf{Pitfall}: Not checking all 6 pairs systematically
    \begin{itemize}
        \item \textbf{Impact}: HIGH - Will miss intersections and get wrong count
        \item \textbf{Prevention}: List all 6 pairs explicitly before starting
        \item \textbf{Detection}: If answer seems too low, recheck which pairs were solved
    \end{itemize}
    
    \item \textbf{Pitfall}: Double-counting points that appear in multiple pairs
    \begin{itemize}
        \item \textbf{Impact}: MEDIUM - Could count 6 or 7 instead of 5
        \item \textbf{Prevention}: Keep organized list of distinct $(x, y)$ coordinates
        \item \textbf{Detection}: Final list should have 5 unique points with different coordinates
    \end{itemize}
    
    \item \textbf{Pitfall}: Sign errors when applying $\log_{\frac{1}{a}} b = -\log_a b$
    \begin{itemize}
        \item \textbf{Impact}: MEDIUM - Could get wrong equations for $f_3$ and $f_4$
        \item \textbf{Prevention}: Write out property carefully and test with specific value
        \item \textbf{Detection}: Verify by computing $\log_{\frac{1}{3}} 3 = -1$ directly
    \end{itemize}
    
    \item \textbf{Pitfall}: Algebraic error when solving $t^2 = 1$ or similar
    \begin{itemize}
        \item \textbf{Impact}: MEDIUM - Missing $t = -1$ solution, only finding $t = 1$
        \item \textbf{Prevention}: Always consider both positive and negative square roots
        \item \textbf{Detection}: Check that number of solutions seems reasonable
    \end{itemize}
    
    \item \textbf{Pitfall}: Forgetting to convert back from $t$-coordinates to $x$-coordinates
    \begin{itemize}
        \item \textbf{Impact}: LOW - Problem asks for points in $(x, y)$ plane
        \item \textbf{Prevention}: Use $x = 3^t$ to convert after solving in $t$-space
        \item \textbf{Detection}: Final answer should have $x$-values like $3, \frac{1}{3}$, not $t$-values like $1, -1$
    \end{itemize}
    
    \item \textbf{Pitfall}: Claiming no intersection exists when algebra gives no solution
    \begin{itemize}
        \item \textbf{Impact}: LOW - This is actually correct behavior for some pairs
        \item \textbf{Prevention}: Don't assume every pair must intersect
        \item \textbf{Detection}: Verify graphically that functions don't intersect
    \end{itemize}
    
    \item \textbf{Pitfall}: Time management - spending too long on verification
    \begin{itemize}
        \item \textbf{Impact}: LOW-MEDIUM - Could lose time for other problems
        \item \textbf{Prevention}: Set 6-minute time limit for this problem
        \item \textbf{Detection}: Check watch after finding 4-5 intersections
    \end{itemize}
\end{enumerate}
\end{roadmapbox}

\section{COMPETITION STRATEGY INTEGRATION}

\begin{competitionbox}{A. Time Management}
\begin{itemize}[leftmargin=*]
    \item \textbf{Target Time}: 5-6 minutes for complete solution
    \begin{itemize}
        \item Initial analysis: 1 minute
        \item Substitution and setup: 1 minute
        \item Solving 6 pairs: 2-3 minutes
        \item Verification: 1 minute
    \end{itemize}
    
    \item \textbf{Time Allocation Strategy}:
    \begin{itemize}
        \item If solution method clear: work steadily through all pairs
        \item If stuck after 3 minutes: pivot to alternative approach or sketch
        \item If approaching 7 minutes: make best guess from partial work and move on
    \end{itemize}
    
    \item \textbf{Opportunity Cost}:
    \begin{itemize}
        \item P16 is moderately difficult, worth same points as easier problems
        \item Don't spend $>$ 8 minutes unless all easier problems completed
        \item Consider skipping if weak on logarithms
    \end{itemize}
\end{itemize}
\end{competitionbox}

\begin{competitionbox}{B. Problem Prioritization}
\begin{itemize}[leftmargin=*]
    \item \textbf{When to Attempt}:
    \begin{itemize}
        \item After completing P1-P15 (or at least P1-P14)
        \item If comfortable with logarithmic properties
        \item If seeking challenge and have time buffer
    \end{itemize}
    
    \item \textbf{When to Skip}:
    \begin{itemize}
        \item If logarithms are weak area
        \item If less than 15 minutes remain and P17-P20 not attempted
        \item If current pace suggests won't finish all problems
    \end{itemize}
    
    \item \textbf{Strategic Positioning}:
    \begin{itemize}
        \item Good problem to attempt if strong in algebra
        \item Systematic approach minimizes risk
        \item Multiple-choice format allows educated guessing if needed
    \end{itemize}
\end{itemize}
\end{competitionbox}

\begin{competitionbox}{C. Risk Management}
\begin{itemize}[leftmargin=*]
    \item \textbf{Risk Assessment}: MEDIUM
    \begin{itemize}
        \item Systematic enumeration reduces risk of missing cases
        \item Clear answer choices provide validation
        \item Domain restrictions create some pitfall risk
    \end{itemize}
    
    \item \textbf{Risk Mitigation}:
    \begin{itemize}
        \item Use substitution method to simplify algebra
        \item List all 6 pairs explicitly before solving
        \item Quick verification by substitution
        \item If answer doesn't match option, signals error
    \end{itemize}
    
    \item \textbf{Confidence Threshold}:
    \begin{itemize}
        \item If $>$ 80\% confident after solution: bubble answer
        \item If 50-80\% confident: quick verification, then bubble
        \item If $<$ 50\% confident: consider leaving blank or strategic guess
    \end{itemize}
\end{itemize}
\end{competitionbox}

\begin{competitionbox}{D. Abandonment Criteria}
\begin{itemize}[leftmargin=*]
    \item \textbf{Hard Abandonment} (must move on):
    \begin{itemize}
        \item After 8 minutes with no clear path to answer
        \item If completely stuck on logarithm properties
        \item If time pressure critical ($<$ 10 minutes remain, many problems left)
    \end{itemize}
    
    \item \textbf{Soft Abandonment} (make strategic guess):
    \begin{itemize}
        \item After 6-7 minutes if only partial progress
        \item If found 3-4 intersections but struggling to find more
        \item Use partial work to eliminate unlikely options
    \end{itemize}
    
    \item \textbf{Return Strategy}:
    \begin{itemize}
        \item Mark problem for return if time permits
        \item Priority for return: MEDIUM (after P17-P25 if easier)
        \item If returning, focus on finding the intersection at $x = 1$
    \end{itemize}
\end{itemize}
\end{competitionbox}

\newpage

\section{SPECIAL HIGHLIGHT BOXES}

\begin{insightbox}
\textbf{The Game-Changing Realization:}

The key insight is recognizing that the substitution $t = \log_3 x$ transforms all four logarithmic functions into simple algebraic expressions:
\begin{align*}
y &= t \quad \text{(linear)} \\
y &= \frac{1}{t} \quad \text{(reciprocal)} \\
y &= -t \quad \text{(reflection of linear)} \\
y &= -\frac{1}{t} \quad \text{(reflection of reciprocal)}
\end{align*}

This transformation reveals the underlying symmetry structure and reduces the problem from transcendental equations to simple polynomial equations. Once in this form, intersection-finding becomes straightforward: $t = \frac{1}{t}$ gives $t^2 = 1$, etc.

\textbf{Why This Works:} All four functions are fundamentally related through the change-of-base formula and reciprocal base property. The substitution makes these relationships explicit and algebraically tractable.
\end{insightbox}

\begin{warningbox}
\textbf{The Most Common Mistake:}

Students frequently \textbf{forget to check whether $x = 1$ is in the domain} of all functions. While $f_1(1) = f_3(1) = 0$ correctly, $f_2$ and $f_4$ are \textbf{undefined} at $x = 1$ because they have $x$ as the base (and $\log_1$ anything is undefined).

\textbf{Critical Check:} When finding the intersection of $f_1$ and $f_3$ at $(1, 0)$, verify that this point is:
\begin{itemize}
    \item $\checkmark$ Valid for $f_1$ and $f_3$ (they do intersect here)
    \item $\times$ Not relevant for checking "two or more graphs" since $f_2$ and $f_4$ don't exist at $x = 1$
\end{itemize}

However, the problem asks for points on "two or more of the graphs," so $(1, 0)$ \textbf{does count} as it's on exactly two graphs ($f_1$ and $f_3$), even though $f_2$ and $f_4$ aren't defined there.

\textbf{Second Warning:} Don't double-count! The point $(3, 1)$ appears when solving $(f_1, f_2)$, but make sure you're not counting it again elsewhere.
\end{warningbox}

\begin{protipbox}
\textbf{Competition Time-Saver:}

\textbf{Technique:} Before solving all 6 pairs algebraically, use \textbf{quick symmetry observations}:
\begin{itemize}
    \item $f_1$ and $f_2$ are mutual reflections across $y = x$
    \item $f_3$ and $f_4$ are also mutual reflections across $y = x$
    \item $f_1$ and $f_3$ are reflections across the $x$-axis (negatives)
    \item $f_2$ and $f_4$ are reflections across the $x$-axis (negatives)
\end{itemize}

\textbf{Implication:} Pairs that are pure negatives of each other ($f_1$ vs $f_3$, $f_2$ vs $f_4$) can only intersect where $y = 0$ or where both are undefined. This immediately tells you:
\begin{itemize}
    \item $(f_1, f_3)$ intersection: only at $x = 1$ where $y = 0$
    \item $(f_2, f_4)$ intersection: only where $\frac{1}{t} = -\frac{1}{t}$, which has no solution
\end{itemize}

This reduces your workload from 6 pairs to focusing on $(f_1, f_2)$, $(f_1, f_4)$, $(f_2, f_3)$, and $(f_3, f_4)$ for non-trivial intersections.

\textbf{Time Saved:} 1-2 minutes by using symmetry to predict which pairs will/won't intersect.
\end{protipbox}

\section{COMPLETE SOLUTION SUMMARY}

\begin{tcolorbox}[colback=gray!10!white, colframe=black, title=Final Answer with Justification]
\textbf{Answer: (D) 5 points}

\textbf{The five intersection points are:}
\begin{enumerate}
    \item $(3, 1)$ - intersection of $f_1$ and $f_2$
    \item $(\frac{1}{3}, -1)$ - intersection of $f_1$ and $f_2$
    \item $(1, 0)$ - intersection of $f_1$ and $f_3$
    \item $(3, -1)$ - intersection of $f_3$ and $f_4$
    \item $(\frac{1}{3}, 1)$ - intersection of $f_3$ and $f_4$
\end{enumerate}

\textbf{Key Solution Steps:}
\begin{itemize}
    \item Applied logarithmic properties to express all functions in terms of $\log_3 x$
    \item Used substitution $t = \log_3 x$ to transform to algebraic equations
    \item Systematically solved all 6 pairwise intersection equations
    \item Verified domain restrictions and eliminated invalid solutions
    \item Confirmed no double-counting of intersection points
\end{itemize}

\textbf{Verification:} Each point verified by direct substitution. All 6 pairs checked. Answer matches option (D).

\textbf{Confidence Level:} 95\%
\end{tcolorbox}

\end{document}
